%% ============================================================
%% Lecture 17: Limited Dependent Variable Models and
%%             Sample Selection Corrections
%% Intermediate Econometrics — SUFE
%% Wooldridge, Introductory Econometrics, 8th edition, Chapter 17
%% ============================================================
\documentclass[aspectratio=169, 12pt]{beamer}
% ==============================================================================
% header.tex — LaTeX Preamble for Intermediate Econometrics
% Shandong University of Finance and Economics (SUFE)
% Textbook: Wooldridge, Introductory Econometrics, 8th edition
% ==============================================================================
%
% USAGE (from Slides/ directory):
%   % ==============================================================================
% header.tex — LaTeX Preamble for Intermediate Econometrics
% Shandong University of Finance and Economics (SUFE)
% Textbook: Wooldridge, Introductory Econometrics, 8th edition
% ==============================================================================
%
% USAGE (from Slides/ directory):
%   % ==============================================================================
% header.tex — LaTeX Preamble for Intermediate Econometrics
% Shandong University of Finance and Economics (SUFE)
% Textbook: Wooldridge, Introductory Econometrics, 8th edition
% ==============================================================================
%
% USAGE (from Slides/ directory):
%   \input{header}          % in Beamer .tex files (TEXINPUTS=../Preambles:)
%
% COMPILE (3-pass XeLaTeX):
%   cd Slides
%   TEXINPUTS=../Preambles:$TEXINPUTS xelatex -interaction=nonstopmode file.tex
%   BIBINPUTS=..:$BIBINPUTS bibtex file
%   TEXINPUTS=../Preambles:$TEXINPUTS xelatex -interaction=nonstopmode file.tex
%   TEXINPUTS=../Preambles:$TEXINPUTS xelatex -interaction=nonstopmode file.tex
% ==============================================================================


% --- Essential packages -------------------------------------------------------
\usepackage{fontspec}           % XeLaTeX font support (must load early)
\usepackage{amsmath}            % Math environments (align, equation, etc.)
\usepackage{amssymb}            % Math symbols (mathbb, etc.)
\usepackage{bm}                 % Bold math (\bm{})
\usepackage{mathtools}          % Extended math tools (coloneqq, etc.)


% --- Table packages -----------------------------------------------------------
\usepackage{booktabs}           % Professional tables (\toprule, \midrule, \bottomrule)
\usepackage{threeparttable}     % Table notes (tablenotes environment)
\usepackage{siunitx}            % Number alignment in tables (S column type)
\sisetup{
  table-format = 1.4,
  table-number-alignment = center,
  detect-weight = true,
  detect-inline-weight = math
}


% --- Figure and graphics packages --------------------------------------------
\usepackage{graphicx}           % Include figures (\includegraphics)
\usepackage{subcaption}         % Subfigures (subfigure environment)
\usepackage{tikz}               % TikZ diagrams
\usepackage{pgfplots}           % Data plots
\pgfplotsset{compat=1.18}
\usetikzlibrary{arrows.meta, positioning, shapes.geometric,
                decorations.pathreplacing, calc, patterns}


% --- Color and boxes ----------------------------------------------------------
\usepackage{xcolor}             % Extended color support
\usepackage{tcolorbox}          % Custom box environments
\tcbuselibrary{skins, breakable, theorems}


% --- Other utility packages ---------------------------------------------------
\usepackage{hyperref}           % Hyperlinks
\usepackage{natbib}             % Bibliography (\citet, \citep, \citealt)
\usepackage{caption}            % Figure/table captions
\usepackage{appendixnumberbeamer} % Appendix frame numbering


% ==============================================================================
% SUFE INSTITUTIONAL COLORS
% ==============================================================================
% Update hex values to match official SUFE brand colors if available.

\definecolor{SUFEBlue}{HTML}{012169}     % Primary: deep navy
\definecolor{SUFEGold}{HTML}{B9975B}     % Secondary: warm gold
\definecolor{SUFEYellow}{HTML}{F2A900}   % Highlight: bright gold/yellow
\definecolor{SUFELight}{HTML}{E8EDF5}    % Light background: pale blue-gray
\definecolor{SUFEDark}{HTML}{1a1a1a}     % Body text: near-black
\definecolor{positive}{HTML}{2E7D32}     % Positive/correct: dark green
\definecolor{negative}{HTML}{C62828}     % Negative/incorrect: dark red


% ==============================================================================
% BEAMER THEME & COLOR SETUP
% ==============================================================================

\usetheme{default}
\useinnertheme{default}
\useoutertheme{default}
\usecolortheme{default}

% Frame title
\setbeamercolor{frametitle}{fg=SUFEBlue, bg=white}
\setbeamerfont{frametitle}{size=\large, series=\bfseries}
\setbeamertemplate{frametitle}[default][left]
\addtobeamertemplate{frametitle}{\vspace{0.05em}}{%
  \vspace{0.05em}{\color{SUFEGold}\rule{\textwidth}{0.5pt}}\vspace{-0.1em}}

% Title page
\setbeamercolor{title}{fg=SUFEBlue}
\setbeamercolor{subtitle}{fg=SUFEGold}
\setbeamercolor{author}{fg=SUFEBlue}
\setbeamercolor{institute}{fg=SUFEDark}
\setbeamercolor{date}{fg=SUFEGold}

% Structure elements
\setbeamercolor{structure}{fg=SUFEBlue}
\setbeamercolor{item}{fg=SUFEBlue}
\setbeamercolor{subitem}{fg=SUFEGold}
\setbeamercolor{subsubitem}{fg=SUFEGold}
\setbeamerfont{itemize/enumerate body}{size=\normalsize}
\setbeamerfont{itemize/enumerate subbody}{size=\small}

% Block environments
\setbeamercolor{block title}{fg=white, bg=SUFEBlue}
\setbeamercolor{block body}{fg=SUFEDark, bg=SUFELight}
\setbeamercolor{block title alerted}{fg=white, bg=red!70!black}
\setbeamercolor{block body alerted}{fg=SUFEDark, bg=red!5}
\setbeamercolor{block title example}{fg=white, bg=SUFEGold!80!black}
\setbeamercolor{block body example}{fg=SUFEDark, bg=SUFEGold!8}

% Remove navigation symbols
\setbeamertemplate{navigation symbols}{}

% Footline: author | short title | slide N/N
\setbeamertemplate{footline}{%
  \leavevmode%
  \hbox{%
    \begin{beamercolorbox}[wd=.30\paperwidth,ht=2.25ex,dp=1ex,left,leftskip=2mm]{author in head/foot}%
      \usebeamerfont{author in head/foot}\insertshortauthor
    \end{beamercolorbox}%
    \begin{beamercolorbox}[wd=.40\paperwidth,ht=2.25ex,dp=1ex,center]{title in head/foot}%
      \usebeamerfont{title in head/foot}\insertshorttitle
    \end{beamercolorbox}%
    \begin{beamercolorbox}[wd=.30\paperwidth,ht=2.25ex,dp=1ex,right,rightskip=2mm]{date in head/foot}%
      \usebeamerfont{date in head/foot}%
      \insertframenumber{} / \inserttotalframenumber
    \end{beamercolorbox}}%
  \vskip0pt%
}
\setbeamercolor{author in head/foot}{fg=white, bg=SUFEBlue}
\setbeamercolor{title in head/foot}{fg=white, bg=SUFEBlue!80!black}
\setbeamercolor{date in head/foot}{fg=white, bg=SUFEBlue}

% Section separator slides (large centered section title)
\AtBeginSection[]{
  \begin{frame}
    \centering
    \vfill
    {\Large\bfseries\color{SUFEBlue}\insertsectionhead}
    \vfill
    {\color{SUFEGold}\rule{0.5\textwidth}{1pt}}
    \vfill
  \end{frame}
}


% ==============================================================================
% FONT SETUP (XeLaTeX)
% ==============================================================================
% Using Windows system fonts (Times New Roman / Arial / Consolas).
% On macOS, swap for: TeX Gyre Termes / TeX Gyre Heros / TeX Gyre Cursor
%   or: Times New Roman / Helvetica / Courier New

\setmainfont{Times New Roman}
\setsansfont{Arial}
\setmonofont{Consolas}


% ==============================================================================
% CUSTOM BOX ENVIRONMENTS (tcolorbox)
% ==============================================================================
% Quarto CSS equivalents: .keybox, .highlightbox, .methodbox,
%                         .assumptionbox, .resultbox

% keybox: key definitions, stated assumptions
\newtcolorbox{keybox}{
  enhanced,
  colback  = SUFEGold!10!white,
  colframe = SUFEGold,
  leftrule = 4pt, rightrule = 0pt, toprule = 0pt, bottomrule = 0pt,
  arc = 0pt, outer arc = 0pt,
  boxsep = 0pt, left = 8pt, right = 6pt, top = 4pt, bottom = 4pt,
}

% highlightbox: important theorems, main results
\newtcolorbox{highlightbox}{
  enhanced,
  colback  = SUFEYellow!8!white,
  colframe = SUFEYellow!80!SUFEGold,
  leftrule = 4pt, rightrule = 0pt, toprule = 0pt, bottomrule = 0pt,
  arc = 0pt, outer arc = 0pt,
  boxsep = 0pt, left = 8pt, right = 6pt, top = 4pt, bottom = 4pt,
}

% definitionbox{Title}: formal definitions — OLS, IV, GMM, etc.
\newtcolorbox{definitionbox}[1]{
  enhanced,
  title    = #1,
  colback  = SUFEBlue!4!white,
  colframe = SUFEBlue,
  colbacktitle = SUFEBlue,
  coltitle = white,
  fonttitle= \bfseries\small,
  arc = 4pt, outer arc = 4pt,
  boxsep = 0pt, left = 8pt, right = 6pt, top = 4pt, bottom = 4pt,
  toptitle = 3pt, bottomtitle = 3pt,
}

% examplebox{Title}: empirical examples, Wooldridge data applications
\newtcolorbox{examplebox}[1]{
  enhanced,
  title    = #1,
  colback  = SUFELight,
  colframe = SUFEGold!70!white,
  colbacktitle = SUFEGold!35!white,
  coltitle = SUFEBlue,
  fonttitle= \bfseries\small,
  arc = 4pt, outer arc = 4pt,
  boxsep = 0pt, left = 8pt, right = 6pt, top = 4pt, bottom = 4pt,
  toptitle = 3pt, bottomtitle = 3pt,
}

% assumptionbox: numbered OLS assumptions (MLR.1–MLR.6)
\newtcolorbox{assumptionbox}{
  enhanced,
  colback  = SUFEGold!8!white,
  colframe = SUFEGold,
  arc = 4pt, outer arc = 4pt,
  boxsep = 0pt, left = 8pt, right = 6pt, top = 4pt, bottom = 4pt,
}


% ==============================================================================
% STANDARD MATH MACROS
% ==============================================================================

% --- Operators ----------------------------------------------------------------
\DeclareMathOperator{\E}{\mathbb{E}}       % Expectation
\DeclareMathOperator{\Var}{Var}             % Variance
\DeclareMathOperator{\Cov}{Cov}             % Covariance
\DeclareMathOperator{\Corr}{Corr}           % Correlation
\DeclareMathOperator{\plim}{plim}           % Probability limit
\DeclareMathOperator{\rank}{rank}           % Matrix rank
\DeclareMathOperator{\tr}{tr}               % Matrix trace
\DeclareMathOperator{\se}{se}               % Standard error

% --- Common shortcuts ---------------------------------------------------------
\newcommand{\bhat}{\hat{\beta}}             % OLS point estimator
\newcommand{\bhatt}{\tilde{\beta}}          % Alternative estimator
\newcommand{\eps}{\varepsilon}              % Error term
\newcommand{\epsh}{\hat{\varepsilon}}       % OLS residual
\newcommand{\uh}{\hat{u}}                   % Residual (alternative)
\newcommand{\yh}{\hat{y}}                   % Fitted value

% Matrices and vectors
\newcommand{\Xb}{\mathbf{X}}               % Design matrix
\newcommand{\yb}{\mathbf{y}}               % Outcome vector
\newcommand{\betab}{\boldsymbol{\beta}}    % Parameter vector
\newcommand{\epsb}{\boldsymbol{\varepsilon}} % Error vector

% Goodness-of-fit
\newcommand{\SSR}{\text{SSR}}              % Sum of squared residuals
\newcommand{\SSE}{\text{SSE}}              % Explained sum of squares
\newcommand{\SST}{\text{SST}}              % Total sum of squares
\newcommand{\Rsq}{R^2}                     % R-squared
\newcommand{\aRsq}{\bar{R}^2}             % Adjusted R-squared

% Asymptotic notation
\newcommand{\avar}{\text{Avar}}            % Asymptotic variance
\newcommand{\pto}{\overset{p}{\to}}        % Convergence in probability
\newcommand{\dto}{\overset{d}{\to}}        % Convergence in distribution
\newcommand{\iid}{\overset{\text{iid}}{\sim}}  % i.i.d. distributed
\newcommand{\asim}{\overset{a}{\sim}}      % Approximately distributed

% Econometric estimators
\newcommand{\OLS}{\text{OLS}}
\newcommand{\IV}{\text{IV}}
\newcommand{\TSLS}{\text{2SLS}}
\newcommand{\GMM}{\text{GMM}}
\newcommand{\FE}{\text{FE}}
\newcommand{\RE}{\text{RE}}
\newcommand{\DID}{\text{DiD}}
\newcommand{\RD}{\text{RD}}

% Probability
\newcommand{\prob}{\mathbb{P}}             % Probability
\newcommand{\indep}{\perp\!\!\!\perp}      % Statistical independence

% Distributions
\newcommand{\Normal}{\mathcal{N}}          % Normal distribution
% Usage: X \sim \Normal(\mu, \sigma^2)

% Absolute value and norm
\newcommand{\abs}[1]{\left\lvert #1 \right\rvert}
\newcommand{\norm}[1]{\left\lVert #1 \right\rVert}


% ==============================================================================
% HYPERREF SETUP
% ==============================================================================

\hypersetup{
  colorlinks = true,
  linkcolor  = SUFEBlue,
  citecolor  = SUFEGold,
  urlcolor   = SUFEBlue
}


% ==============================================================================
% BIBLIOGRAPHY SETUP
% ==============================================================================

\bibliographystyle{aer}   % American Economic Review style
% Usage: \bibliography{../Bibliography_base}  (from Slides/ directory)
% In-text: \citet{Wooldridge2020}, \citep{White1980}
          % in Beamer .tex files (TEXINPUTS=../Preambles:)
%
% COMPILE (3-pass XeLaTeX):
%   cd Slides
%   TEXINPUTS=../Preambles:$TEXINPUTS xelatex -interaction=nonstopmode file.tex
%   BIBINPUTS=..:$BIBINPUTS bibtex file
%   TEXINPUTS=../Preambles:$TEXINPUTS xelatex -interaction=nonstopmode file.tex
%   TEXINPUTS=../Preambles:$TEXINPUTS xelatex -interaction=nonstopmode file.tex
% ==============================================================================


% --- Essential packages -------------------------------------------------------
\usepackage{fontspec}           % XeLaTeX font support (must load early)
\usepackage{amsmath}            % Math environments (align, equation, etc.)
\usepackage{amssymb}            % Math symbols (mathbb, etc.)
\usepackage{bm}                 % Bold math (\bm{})
\usepackage{mathtools}          % Extended math tools (coloneqq, etc.)


% --- Table packages -----------------------------------------------------------
\usepackage{booktabs}           % Professional tables (\toprule, \midrule, \bottomrule)
\usepackage{threeparttable}     % Table notes (tablenotes environment)
\usepackage{siunitx}            % Number alignment in tables (S column type)
\sisetup{
  table-format = 1.4,
  table-number-alignment = center,
  detect-weight = true,
  detect-inline-weight = math
}


% --- Figure and graphics packages --------------------------------------------
\usepackage{graphicx}           % Include figures (\includegraphics)
\usepackage{subcaption}         % Subfigures (subfigure environment)
\usepackage{tikz}               % TikZ diagrams
\usepackage{pgfplots}           % Data plots
\pgfplotsset{compat=1.18}
\usetikzlibrary{arrows.meta, positioning, shapes.geometric,
                decorations.pathreplacing, calc, patterns}


% --- Color and boxes ----------------------------------------------------------
\usepackage{xcolor}             % Extended color support
\usepackage{tcolorbox}          % Custom box environments
\tcbuselibrary{skins, breakable, theorems}


% --- Other utility packages ---------------------------------------------------
\usepackage{hyperref}           % Hyperlinks
\usepackage{natbib}             % Bibliography (\citet, \citep, \citealt)
\usepackage{caption}            % Figure/table captions
\usepackage{appendixnumberbeamer} % Appendix frame numbering


% ==============================================================================
% SUFE INSTITUTIONAL COLORS
% ==============================================================================
% Update hex values to match official SUFE brand colors if available.

\definecolor{SUFEBlue}{HTML}{012169}     % Primary: deep navy
\definecolor{SUFEGold}{HTML}{B9975B}     % Secondary: warm gold
\definecolor{SUFEYellow}{HTML}{F2A900}   % Highlight: bright gold/yellow
\definecolor{SUFELight}{HTML}{E8EDF5}    % Light background: pale blue-gray
\definecolor{SUFEDark}{HTML}{1a1a1a}     % Body text: near-black
\definecolor{positive}{HTML}{2E7D32}     % Positive/correct: dark green
\definecolor{negative}{HTML}{C62828}     % Negative/incorrect: dark red


% ==============================================================================
% BEAMER THEME & COLOR SETUP
% ==============================================================================

\usetheme{default}
\useinnertheme{default}
\useoutertheme{default}
\usecolortheme{default}

% Frame title
\setbeamercolor{frametitle}{fg=SUFEBlue, bg=white}
\setbeamerfont{frametitle}{size=\large, series=\bfseries}
\setbeamertemplate{frametitle}[default][left]
\addtobeamertemplate{frametitle}{\vspace{0.05em}}{%
  \vspace{0.05em}{\color{SUFEGold}\rule{\textwidth}{0.5pt}}\vspace{-0.1em}}

% Title page
\setbeamercolor{title}{fg=SUFEBlue}
\setbeamercolor{subtitle}{fg=SUFEGold}
\setbeamercolor{author}{fg=SUFEBlue}
\setbeamercolor{institute}{fg=SUFEDark}
\setbeamercolor{date}{fg=SUFEGold}

% Structure elements
\setbeamercolor{structure}{fg=SUFEBlue}
\setbeamercolor{item}{fg=SUFEBlue}
\setbeamercolor{subitem}{fg=SUFEGold}
\setbeamercolor{subsubitem}{fg=SUFEGold}
\setbeamerfont{itemize/enumerate body}{size=\normalsize}
\setbeamerfont{itemize/enumerate subbody}{size=\small}

% Block environments
\setbeamercolor{block title}{fg=white, bg=SUFEBlue}
\setbeamercolor{block body}{fg=SUFEDark, bg=SUFELight}
\setbeamercolor{block title alerted}{fg=white, bg=red!70!black}
\setbeamercolor{block body alerted}{fg=SUFEDark, bg=red!5}
\setbeamercolor{block title example}{fg=white, bg=SUFEGold!80!black}
\setbeamercolor{block body example}{fg=SUFEDark, bg=SUFEGold!8}

% Remove navigation symbols
\setbeamertemplate{navigation symbols}{}

% Footline: author | short title | slide N/N
\setbeamertemplate{footline}{%
  \leavevmode%
  \hbox{%
    \begin{beamercolorbox}[wd=.30\paperwidth,ht=2.25ex,dp=1ex,left,leftskip=2mm]{author in head/foot}%
      \usebeamerfont{author in head/foot}\insertshortauthor
    \end{beamercolorbox}%
    \begin{beamercolorbox}[wd=.40\paperwidth,ht=2.25ex,dp=1ex,center]{title in head/foot}%
      \usebeamerfont{title in head/foot}\insertshorttitle
    \end{beamercolorbox}%
    \begin{beamercolorbox}[wd=.30\paperwidth,ht=2.25ex,dp=1ex,right,rightskip=2mm]{date in head/foot}%
      \usebeamerfont{date in head/foot}%
      \insertframenumber{} / \inserttotalframenumber
    \end{beamercolorbox}}%
  \vskip0pt%
}
\setbeamercolor{author in head/foot}{fg=white, bg=SUFEBlue}
\setbeamercolor{title in head/foot}{fg=white, bg=SUFEBlue!80!black}
\setbeamercolor{date in head/foot}{fg=white, bg=SUFEBlue}

% Section separator slides (large centered section title)
\AtBeginSection[]{
  \begin{frame}
    \centering
    \vfill
    {\Large\bfseries\color{SUFEBlue}\insertsectionhead}
    \vfill
    {\color{SUFEGold}\rule{0.5\textwidth}{1pt}}
    \vfill
  \end{frame}
}


% ==============================================================================
% FONT SETUP (XeLaTeX)
% ==============================================================================
% Using Windows system fonts (Times New Roman / Arial / Consolas).
% On macOS, swap for: TeX Gyre Termes / TeX Gyre Heros / TeX Gyre Cursor
%   or: Times New Roman / Helvetica / Courier New

\setmainfont{Times New Roman}
\setsansfont{Arial}
\setmonofont{Consolas}


% ==============================================================================
% CUSTOM BOX ENVIRONMENTS (tcolorbox)
% ==============================================================================
% Quarto CSS equivalents: .keybox, .highlightbox, .methodbox,
%                         .assumptionbox, .resultbox

% keybox: key definitions, stated assumptions
\newtcolorbox{keybox}{
  enhanced,
  colback  = SUFEGold!10!white,
  colframe = SUFEGold,
  leftrule = 4pt, rightrule = 0pt, toprule = 0pt, bottomrule = 0pt,
  arc = 0pt, outer arc = 0pt,
  boxsep = 0pt, left = 8pt, right = 6pt, top = 4pt, bottom = 4pt,
}

% highlightbox: important theorems, main results
\newtcolorbox{highlightbox}{
  enhanced,
  colback  = SUFEYellow!8!white,
  colframe = SUFEYellow!80!SUFEGold,
  leftrule = 4pt, rightrule = 0pt, toprule = 0pt, bottomrule = 0pt,
  arc = 0pt, outer arc = 0pt,
  boxsep = 0pt, left = 8pt, right = 6pt, top = 4pt, bottom = 4pt,
}

% definitionbox{Title}: formal definitions — OLS, IV, GMM, etc.
\newtcolorbox{definitionbox}[1]{
  enhanced,
  title    = #1,
  colback  = SUFEBlue!4!white,
  colframe = SUFEBlue,
  colbacktitle = SUFEBlue,
  coltitle = white,
  fonttitle= \bfseries\small,
  arc = 4pt, outer arc = 4pt,
  boxsep = 0pt, left = 8pt, right = 6pt, top = 4pt, bottom = 4pt,
  toptitle = 3pt, bottomtitle = 3pt,
}

% examplebox{Title}: empirical examples, Wooldridge data applications
\newtcolorbox{examplebox}[1]{
  enhanced,
  title    = #1,
  colback  = SUFELight,
  colframe = SUFEGold!70!white,
  colbacktitle = SUFEGold!35!white,
  coltitle = SUFEBlue,
  fonttitle= \bfseries\small,
  arc = 4pt, outer arc = 4pt,
  boxsep = 0pt, left = 8pt, right = 6pt, top = 4pt, bottom = 4pt,
  toptitle = 3pt, bottomtitle = 3pt,
}

% assumptionbox: numbered OLS assumptions (MLR.1–MLR.6)
\newtcolorbox{assumptionbox}{
  enhanced,
  colback  = SUFEGold!8!white,
  colframe = SUFEGold,
  arc = 4pt, outer arc = 4pt,
  boxsep = 0pt, left = 8pt, right = 6pt, top = 4pt, bottom = 4pt,
}


% ==============================================================================
% STANDARD MATH MACROS
% ==============================================================================

% --- Operators ----------------------------------------------------------------
\DeclareMathOperator{\E}{\mathbb{E}}       % Expectation
\DeclareMathOperator{\Var}{Var}             % Variance
\DeclareMathOperator{\Cov}{Cov}             % Covariance
\DeclareMathOperator{\Corr}{Corr}           % Correlation
\DeclareMathOperator{\plim}{plim}           % Probability limit
\DeclareMathOperator{\rank}{rank}           % Matrix rank
\DeclareMathOperator{\tr}{tr}               % Matrix trace
\DeclareMathOperator{\se}{se}               % Standard error

% --- Common shortcuts ---------------------------------------------------------
\newcommand{\bhat}{\hat{\beta}}             % OLS point estimator
\newcommand{\bhatt}{\tilde{\beta}}          % Alternative estimator
\newcommand{\eps}{\varepsilon}              % Error term
\newcommand{\epsh}{\hat{\varepsilon}}       % OLS residual
\newcommand{\uh}{\hat{u}}                   % Residual (alternative)
\newcommand{\yh}{\hat{y}}                   % Fitted value

% Matrices and vectors
\newcommand{\Xb}{\mathbf{X}}               % Design matrix
\newcommand{\yb}{\mathbf{y}}               % Outcome vector
\newcommand{\betab}{\boldsymbol{\beta}}    % Parameter vector
\newcommand{\epsb}{\boldsymbol{\varepsilon}} % Error vector

% Goodness-of-fit
\newcommand{\SSR}{\text{SSR}}              % Sum of squared residuals
\newcommand{\SSE}{\text{SSE}}              % Explained sum of squares
\newcommand{\SST}{\text{SST}}              % Total sum of squares
\newcommand{\Rsq}{R^2}                     % R-squared
\newcommand{\aRsq}{\bar{R}^2}             % Adjusted R-squared

% Asymptotic notation
\newcommand{\avar}{\text{Avar}}            % Asymptotic variance
\newcommand{\pto}{\overset{p}{\to}}        % Convergence in probability
\newcommand{\dto}{\overset{d}{\to}}        % Convergence in distribution
\newcommand{\iid}{\overset{\text{iid}}{\sim}}  % i.i.d. distributed
\newcommand{\asim}{\overset{a}{\sim}}      % Approximately distributed

% Econometric estimators
\newcommand{\OLS}{\text{OLS}}
\newcommand{\IV}{\text{IV}}
\newcommand{\TSLS}{\text{2SLS}}
\newcommand{\GMM}{\text{GMM}}
\newcommand{\FE}{\text{FE}}
\newcommand{\RE}{\text{RE}}
\newcommand{\DID}{\text{DiD}}
\newcommand{\RD}{\text{RD}}

% Probability
\newcommand{\prob}{\mathbb{P}}             % Probability
\newcommand{\indep}{\perp\!\!\!\perp}      % Statistical independence

% Distributions
\newcommand{\Normal}{\mathcal{N}}          % Normal distribution
% Usage: X \sim \Normal(\mu, \sigma^2)

% Absolute value and norm
\newcommand{\abs}[1]{\left\lvert #1 \right\rvert}
\newcommand{\norm}[1]{\left\lVert #1 \right\rVert}


% ==============================================================================
% HYPERREF SETUP
% ==============================================================================

\hypersetup{
  colorlinks = true,
  linkcolor  = SUFEBlue,
  citecolor  = SUFEGold,
  urlcolor   = SUFEBlue
}


% ==============================================================================
% BIBLIOGRAPHY SETUP
% ==============================================================================

\bibliographystyle{aer}   % American Economic Review style
% Usage: \bibliography{../Bibliography_base}  (from Slides/ directory)
% In-text: \citet{Wooldridge2020}, \citep{White1980}
          % in Beamer .tex files (TEXINPUTS=../Preambles:)
%
% COMPILE (3-pass XeLaTeX):
%   cd Slides
%   TEXINPUTS=../Preambles:$TEXINPUTS xelatex -interaction=nonstopmode file.tex
%   BIBINPUTS=..:$BIBINPUTS bibtex file
%   TEXINPUTS=../Preambles:$TEXINPUTS xelatex -interaction=nonstopmode file.tex
%   TEXINPUTS=../Preambles:$TEXINPUTS xelatex -interaction=nonstopmode file.tex
% ==============================================================================


% --- Essential packages -------------------------------------------------------
\usepackage{fontspec}           % XeLaTeX font support (must load early)
\usepackage{amsmath}            % Math environments (align, equation, etc.)
\usepackage{amssymb}            % Math symbols (mathbb, etc.)
\usepackage{bm}                 % Bold math (\bm{})
\usepackage{mathtools}          % Extended math tools (coloneqq, etc.)


% --- Table packages -----------------------------------------------------------
\usepackage{booktabs}           % Professional tables (\toprule, \midrule, \bottomrule)
\usepackage{threeparttable}     % Table notes (tablenotes environment)
\usepackage{siunitx}            % Number alignment in tables (S column type)
\sisetup{
  table-format = 1.4,
  table-number-alignment = center,
  detect-weight = true,
  detect-inline-weight = math
}


% --- Figure and graphics packages --------------------------------------------
\usepackage{graphicx}           % Include figures (\includegraphics)
\usepackage{subcaption}         % Subfigures (subfigure environment)
\usepackage{tikz}               % TikZ diagrams
\usepackage{pgfplots}           % Data plots
\pgfplotsset{compat=1.18}
\usetikzlibrary{arrows.meta, positioning, shapes.geometric,
                decorations.pathreplacing, calc, patterns}


% --- Color and boxes ----------------------------------------------------------
\usepackage{xcolor}             % Extended color support
\usepackage{tcolorbox}          % Custom box environments
\tcbuselibrary{skins, breakable, theorems}


% --- Other utility packages ---------------------------------------------------
\usepackage{hyperref}           % Hyperlinks
\usepackage{natbib}             % Bibliography (\citet, \citep, \citealt)
\usepackage{caption}            % Figure/table captions
\usepackage{appendixnumberbeamer} % Appendix frame numbering


% ==============================================================================
% SUFE INSTITUTIONAL COLORS
% ==============================================================================
% Update hex values to match official SUFE brand colors if available.

\definecolor{SUFEBlue}{HTML}{012169}     % Primary: deep navy
\definecolor{SUFEGold}{HTML}{B9975B}     % Secondary: warm gold
\definecolor{SUFEYellow}{HTML}{F2A900}   % Highlight: bright gold/yellow
\definecolor{SUFELight}{HTML}{E8EDF5}    % Light background: pale blue-gray
\definecolor{SUFEDark}{HTML}{1a1a1a}     % Body text: near-black
\definecolor{positive}{HTML}{2E7D32}     % Positive/correct: dark green
\definecolor{negative}{HTML}{C62828}     % Negative/incorrect: dark red


% ==============================================================================
% BEAMER THEME & COLOR SETUP
% ==============================================================================

\usetheme{default}
\useinnertheme{default}
\useoutertheme{default}
\usecolortheme{default}

% Frame title
\setbeamercolor{frametitle}{fg=SUFEBlue, bg=white}
\setbeamerfont{frametitle}{size=\large, series=\bfseries}
\setbeamertemplate{frametitle}[default][left]
\addtobeamertemplate{frametitle}{\vspace{0.05em}}{%
  \vspace{0.05em}{\color{SUFEGold}\rule{\textwidth}{0.5pt}}\vspace{-0.1em}}

% Title page
\setbeamercolor{title}{fg=SUFEBlue}
\setbeamercolor{subtitle}{fg=SUFEGold}
\setbeamercolor{author}{fg=SUFEBlue}
\setbeamercolor{institute}{fg=SUFEDark}
\setbeamercolor{date}{fg=SUFEGold}

% Structure elements
\setbeamercolor{structure}{fg=SUFEBlue}
\setbeamercolor{item}{fg=SUFEBlue}
\setbeamercolor{subitem}{fg=SUFEGold}
\setbeamercolor{subsubitem}{fg=SUFEGold}
\setbeamerfont{itemize/enumerate body}{size=\normalsize}
\setbeamerfont{itemize/enumerate subbody}{size=\small}

% Block environments
\setbeamercolor{block title}{fg=white, bg=SUFEBlue}
\setbeamercolor{block body}{fg=SUFEDark, bg=SUFELight}
\setbeamercolor{block title alerted}{fg=white, bg=red!70!black}
\setbeamercolor{block body alerted}{fg=SUFEDark, bg=red!5}
\setbeamercolor{block title example}{fg=white, bg=SUFEGold!80!black}
\setbeamercolor{block body example}{fg=SUFEDark, bg=SUFEGold!8}

% Remove navigation symbols
\setbeamertemplate{navigation symbols}{}

% Footline: author | short title | slide N/N
\setbeamertemplate{footline}{%
  \leavevmode%
  \hbox{%
    \begin{beamercolorbox}[wd=.30\paperwidth,ht=2.25ex,dp=1ex,left,leftskip=2mm]{author in head/foot}%
      \usebeamerfont{author in head/foot}\insertshortauthor
    \end{beamercolorbox}%
    \begin{beamercolorbox}[wd=.40\paperwidth,ht=2.25ex,dp=1ex,center]{title in head/foot}%
      \usebeamerfont{title in head/foot}\insertshorttitle
    \end{beamercolorbox}%
    \begin{beamercolorbox}[wd=.30\paperwidth,ht=2.25ex,dp=1ex,right,rightskip=2mm]{date in head/foot}%
      \usebeamerfont{date in head/foot}%
      \insertframenumber{} / \inserttotalframenumber
    \end{beamercolorbox}}%
  \vskip0pt%
}
\setbeamercolor{author in head/foot}{fg=white, bg=SUFEBlue}
\setbeamercolor{title in head/foot}{fg=white, bg=SUFEBlue!80!black}
\setbeamercolor{date in head/foot}{fg=white, bg=SUFEBlue}

% Section separator slides (large centered section title)
\AtBeginSection[]{
  \begin{frame}
    \centering
    \vfill
    {\Large\bfseries\color{SUFEBlue}\insertsectionhead}
    \vfill
    {\color{SUFEGold}\rule{0.5\textwidth}{1pt}}
    \vfill
  \end{frame}
}


% ==============================================================================
% FONT SETUP (XeLaTeX)
% ==============================================================================
% Using Windows system fonts (Times New Roman / Arial / Consolas).
% On macOS, swap for: TeX Gyre Termes / TeX Gyre Heros / TeX Gyre Cursor
%   or: Times New Roman / Helvetica / Courier New

\setmainfont{Times New Roman}
\setsansfont{Arial}
\setmonofont{Consolas}


% ==============================================================================
% CUSTOM BOX ENVIRONMENTS (tcolorbox)
% ==============================================================================
% Quarto CSS equivalents: .keybox, .highlightbox, .methodbox,
%                         .assumptionbox, .resultbox

% keybox: key definitions, stated assumptions
\newtcolorbox{keybox}{
  enhanced,
  colback  = SUFEGold!10!white,
  colframe = SUFEGold,
  leftrule = 4pt, rightrule = 0pt, toprule = 0pt, bottomrule = 0pt,
  arc = 0pt, outer arc = 0pt,
  boxsep = 0pt, left = 8pt, right = 6pt, top = 4pt, bottom = 4pt,
}

% highlightbox: important theorems, main results
\newtcolorbox{highlightbox}{
  enhanced,
  colback  = SUFEYellow!8!white,
  colframe = SUFEYellow!80!SUFEGold,
  leftrule = 4pt, rightrule = 0pt, toprule = 0pt, bottomrule = 0pt,
  arc = 0pt, outer arc = 0pt,
  boxsep = 0pt, left = 8pt, right = 6pt, top = 4pt, bottom = 4pt,
}

% definitionbox{Title}: formal definitions — OLS, IV, GMM, etc.
\newtcolorbox{definitionbox}[1]{
  enhanced,
  title    = #1,
  colback  = SUFEBlue!4!white,
  colframe = SUFEBlue,
  colbacktitle = SUFEBlue,
  coltitle = white,
  fonttitle= \bfseries\small,
  arc = 4pt, outer arc = 4pt,
  boxsep = 0pt, left = 8pt, right = 6pt, top = 4pt, bottom = 4pt,
  toptitle = 3pt, bottomtitle = 3pt,
}

% examplebox{Title}: empirical examples, Wooldridge data applications
\newtcolorbox{examplebox}[1]{
  enhanced,
  title    = #1,
  colback  = SUFELight,
  colframe = SUFEGold!70!white,
  colbacktitle = SUFEGold!35!white,
  coltitle = SUFEBlue,
  fonttitle= \bfseries\small,
  arc = 4pt, outer arc = 4pt,
  boxsep = 0pt, left = 8pt, right = 6pt, top = 4pt, bottom = 4pt,
  toptitle = 3pt, bottomtitle = 3pt,
}

% assumptionbox: numbered OLS assumptions (MLR.1–MLR.6)
\newtcolorbox{assumptionbox}{
  enhanced,
  colback  = SUFEGold!8!white,
  colframe = SUFEGold,
  arc = 4pt, outer arc = 4pt,
  boxsep = 0pt, left = 8pt, right = 6pt, top = 4pt, bottom = 4pt,
}


% ==============================================================================
% STANDARD MATH MACROS
% ==============================================================================

% --- Operators ----------------------------------------------------------------
\DeclareMathOperator{\E}{\mathbb{E}}       % Expectation
\DeclareMathOperator{\Var}{Var}             % Variance
\DeclareMathOperator{\Cov}{Cov}             % Covariance
\DeclareMathOperator{\Corr}{Corr}           % Correlation
\DeclareMathOperator{\plim}{plim}           % Probability limit
\DeclareMathOperator{\rank}{rank}           % Matrix rank
\DeclareMathOperator{\tr}{tr}               % Matrix trace
\DeclareMathOperator{\se}{se}               % Standard error

% --- Common shortcuts ---------------------------------------------------------
\newcommand{\bhat}{\hat{\beta}}             % OLS point estimator
\newcommand{\bhatt}{\tilde{\beta}}          % Alternative estimator
\newcommand{\eps}{\varepsilon}              % Error term
\newcommand{\epsh}{\hat{\varepsilon}}       % OLS residual
\newcommand{\uh}{\hat{u}}                   % Residual (alternative)
\newcommand{\yh}{\hat{y}}                   % Fitted value

% Matrices and vectors
\newcommand{\Xb}{\mathbf{X}}               % Design matrix
\newcommand{\yb}{\mathbf{y}}               % Outcome vector
\newcommand{\betab}{\boldsymbol{\beta}}    % Parameter vector
\newcommand{\epsb}{\boldsymbol{\varepsilon}} % Error vector

% Goodness-of-fit
\newcommand{\SSR}{\text{SSR}}              % Sum of squared residuals
\newcommand{\SSE}{\text{SSE}}              % Explained sum of squares
\newcommand{\SST}{\text{SST}}              % Total sum of squares
\newcommand{\Rsq}{R^2}                     % R-squared
\newcommand{\aRsq}{\bar{R}^2}             % Adjusted R-squared

% Asymptotic notation
\newcommand{\avar}{\text{Avar}}            % Asymptotic variance
\newcommand{\pto}{\overset{p}{\to}}        % Convergence in probability
\newcommand{\dto}{\overset{d}{\to}}        % Convergence in distribution
\newcommand{\iid}{\overset{\text{iid}}{\sim}}  % i.i.d. distributed
\newcommand{\asim}{\overset{a}{\sim}}      % Approximately distributed

% Econometric estimators
\newcommand{\OLS}{\text{OLS}}
\newcommand{\IV}{\text{IV}}
\newcommand{\TSLS}{\text{2SLS}}
\newcommand{\GMM}{\text{GMM}}
\newcommand{\FE}{\text{FE}}
\newcommand{\RE}{\text{RE}}
\newcommand{\DID}{\text{DiD}}
\newcommand{\RD}{\text{RD}}

% Probability
\newcommand{\prob}{\mathbb{P}}             % Probability
\newcommand{\indep}{\perp\!\!\!\perp}      % Statistical independence

% Distributions
\newcommand{\Normal}{\mathcal{N}}          % Normal distribution
% Usage: X \sim \Normal(\mu, \sigma^2)

% Absolute value and norm
\newcommand{\abs}[1]{\left\lvert #1 \right\rvert}
\newcommand{\norm}[1]{\left\lVert #1 \right\rVert}


% ==============================================================================
% HYPERREF SETUP
% ==============================================================================

\hypersetup{
  colorlinks = true,
  linkcolor  = SUFEBlue,
  citecolor  = SUFEGold,
  urlcolor   = SUFEBlue
}


% ==============================================================================
% BIBLIOGRAPHY SETUP
% ==============================================================================

\bibliographystyle{aer}   % American Economic Review style
% Usage: \bibliography{../Bibliography_base}  (from Slides/ directory)
% In-text: \citet{Wooldridge2020}, \citep{White1980}


%% --- Lecture-local macros (checked against notation registry) ---
\newcommand{\xbeta}{\mathbf{x}\boldsymbol{\beta}}   % linear index x*beta
\newcommand{\xib}{\mathbf{x}_i\boldsymbol{\beta}}    % index for obs i
\newcommand{\xbar}{\bar{\mathbf{x}}}                 % vector of means
\DeclareMathOperator{\APE}{APE}
\DeclareMathOperator{\PEA}{PEA}

\title[L17: Limited Dependent Variables]{Limited Dependent Variable Models \\ and Sample Selection Corrections}
\subtitle{Intermediate Econometrics}
\author[SUFE Econometrics]{Chengye Jia}
\institute{%
  Department of Economics \\
  Shandong University of Finance and Economics}
\date{Spring 2026 \\ \smallskip \scriptsize Wooldridge, Chapter 17}

%% ============================================================
\begin{document}
%% ============================================================

\begin{frame}[plain]
  \titlepage
\end{frame}

\begin{frame}[shrink=10]{Outline}
  \tableofcontents
\end{frame}

%% ============================================================
\section{Motivation: The Limited Dependent Variable Zoo (\S17-0)}
%% ============================================================

\begin{frame}[shrink=8]{Where We Are: From OLS to Nonlinear Models}

  Through Chapter 16 the dependent variable $y$ was treated as (roughly)
  \textbf{continuous and unrestricted}. OLS and its extensions did the work.

  \smallskip

  \begin{keybox}
    \textbf{A limited dependent variable (LDV)} is a variable whose range is
    \textbf{substantively restricted}:
    \begin{itemize}
      \item \textbf{Binary:} $y \in \{0,1\}$ --- employed or not, approved or not
      \item \textbf{Fractional:} $y \in [0,1]$ --- the 401(k) participation rate
      \item \textbf{Count:} $y \in \{0,1,2,\ldots\}$ --- number of arrests in a year
      \item \textbf{Corner solution:} $y \ge 0$ with a pileup at $0$ --- hours worked, charitable gifts
    \end{itemize}
  \end{keybox}

  \smallskip

  \begin{highlightbox}
    \textbf{Why care?} Forcing OLS on these variables can give nonsensical
    fitted values (negative probabilities, negative hours) and constant partial
    effects where economics demands diminishing ones.
  \end{highlightbox}

\end{frame}

\begin{frame}[shrink=8]{Two Distinct Problems Live in This Chapter}

  \begin{columns}[t]
    \column{0.49\textwidth}
    \begin{definitionbox}{Problem 1: The variable itself}
      $y$ is genuinely restricted in the \textbf{population} (observed fully):
      \begin{itemize}
        \item Binary / fractional $\Rightarrow$ logit, probit
        \item Count / nonnegative $\Rightarrow$ Poisson
        \item Corner solution $\Rightarrow$ Tobit
      \end{itemize}
    \end{definitionbox}

    \column{0.49\textwidth}
    \begin{definitionbox}{Problem 2: How we observe $y$}
      $y$ would be fine, but the \textbf{data} are incomplete:
      \begin{itemize}
        \item Top coding / censoring $\Rightarrow$ censored reg.
        \item Sub-population sampled $\Rightarrow$ truncated reg.
        \item Selection on an outcome $\Rightarrow$ Heckit
      \end{itemize}
    \end{definitionbox}
  \end{columns}

  \smallskip

  \begin{keybox}
    \small
    \textbf{Keep these separate.} Tobit for ``hours worked'' models economic
    behavior (people \emph{choose} zero); censored regression for ``top-coded
    wealth'' fixes a survey defect. Same algebra --- very different interpretation.
  \end{keybox}

\end{frame}

\begin{frame}[shrink=8]{Running Application: Married Women's Labor Supply (MROZ)}

  We thread one dataset through the whole lecture --- \texttt{MROZ}, $n=753$
  married women \citep{Mroz1987_sensitivity}.

  \smallskip

  \begin{examplebox}{One dataset --- four LDV faces}
    \begin{itemize}
      \item \textbf{Binary:} $inlf$ --- in the labor force? ($428$ yes, $325$ no)
        $\;\Rightarrow$ logit / probit
      \item \textbf{Corner solution:} $hours$ --- annual hours worked ($325$ zeros,
        positives range $12$--$4{,}950$) $\;\Rightarrow$ Tobit
      \item \textbf{Selection:} $\log(wage)$ --- observed \emph{only} for the $428$
        workers $\;\Rightarrow$ Heckit
      \item \textbf{Fractional / count:} mirrored in companion datasets (401(k)
        $prate$, arrests $narr86$)
    \end{itemize}
  \end{examplebox}

  \smallskip

  \begin{highlightbox}
    Explanatory variables throughout: $nwifeinc$ (other income), $educ$, $exper$,
    $exper^2$, $age$, $kidslt6$ (kids $<6$), $kidsge6$ (kids $6$--$18$).
  \end{highlightbox}

\end{frame}

\begin{frame}[shrink=10]{Socratic Warm-Up}

  \begin{keybox}
    \textbf{Q1.} A linear probability model predicts $\widehat{P}(inlf=1) = 1.08$
    for a highly educated woman. What went wrong, and why can it \emph{never}
    happen with a probit?

    \medskip

    \textbf{Q2.} In the \texttt{MROZ} data, $325$ of $753$ women report exactly
    zero hours. Is ``hours worked'' a censored variable or a corner solution? Does
    the distinction change the model we fit?

    \medskip

    \textbf{Q3.} We only observe wages for women who \emph{work}. If working women
    are systematically more able, will OLS of $\log(wage)$ on $educ$ (using workers
    only) over- or under-state the return to education?
  \end{keybox}

  \smallskip

  We will answer all three by the end of the lecture.

\end{frame}

%% ============================================================
\section{Binary Response: Logit and Probit (\S17-1)}
%% ============================================================

\begin{frame}[shrink=12]{The Linear Probability Model and Its Discontents}

  For a binary $y$, the LPM (Chapter 7) models the \textbf{response probability}
  \[
    P(y=1 \mid \mathbf{x}) = \beta_0 + \beta_1 x_1 + \cdots + \beta_k x_k .
  \]

  \begin{columns}[t]
    \column{0.52\textwidth}
    \begin{highlightbox}
      \textbf{\textcolor{positive}{What LPM gets right}}
      \begin{itemize}
        \item Coefficients \emph{are} partial effects
        \item Consistent for the APE under MLR.1--4
        \item Fast; great for treatment effects
      \end{itemize}
    \end{highlightbox}

    \column{0.46\textwidth}
    \begin{highlightbox}
      \textbf{\textcolor{negative}{What LPM gets wrong}}
      \begin{itemize}
        \item Fitted $\widehat{P}$ can fall outside $[0,1]$
        \item Partial effect of $x_j$ is \emph{constant}
        \item Errors are inherently heteroskedastic
      \end{itemize}
    \end{highlightbox}
  \end{columns}

  \smallskip

  \begin{keybox}
    \textbf{Fix:} run the index $\xbeta$ through a function $G(\cdot)$ that is
    \emph{squashed} into $(0,1)$. That is exactly what logit and probit do.
  \end{keybox}

\end{frame}

\begin{frame}[shrink=12]{Logit and Probit: Squashing the Index into $(0,1)$}

  \begin{definitionbox}{Binary response model}
    \[
      P(y=1 \mid \mathbf{x}) = G(\beta_0 + \beta_1 x_1 + \cdots + \beta_k x_k)
        = G(\beta_0 + \xbeta), \qquad 0 < G(z) < 1 .
    \]
  \end{definitionbox}

  \begin{columns}[c]
    \column{0.36\textwidth}
    \textbf{Logit:} logistic CDF
    \[
      \Lambda(z) = \frac{\exp(z)}{1+\exp(z)} .
    \]
    \textbf{Probit:} normal CDF
    \[
      \Phi(z) = \int_{-\infty}^{z}\!\! \phi(v)\, dv .
    \]

    \column{0.62\textwidth}
    \centering
    \begin{tikzpicture}[scale=0.72]
      \draw[->,thick] (-3.4,0) -- (3.6,0) node[right] {$z$};
      \draw[->,thick] (0,-0.2) -- (0,2.0) node[left] {$G(z)$};
      \draw[dashed,gray!50] (-3.4,1.6) -- (3.4,1.6) node[right,black,font=\scriptsize] {$1$};
      \draw[dashed,gray!50] (-3.4,0.8) -- (3.4,0.8) node[right,black,font=\scriptsize] {$0.5$};
      \draw[thick,SUFEBlue,domain=-3.4:3.4,samples=120,smooth]
        plot (\x, {1.6/(1+exp(-1.6*\x))});
      \draw[thick,SUFEGold,domain=-3.4:3.4,samples=120,smooth]
        plot (\x, {1.6/(1+exp(-2.5*\x))});
      \node[SUFEBlue,font=\scriptsize] at (2.4,1.05) {logit};
      \node[SUFEGold,font=\scriptsize] at (1.1,1.45) {probit};
    \end{tikzpicture}
  \end{columns}

  \begin{keybox}
    \small
    Both are increasing CDFs, steepest at $z=0$. They guarantee
    $\widehat{P}\in(0,1)$ and \emph{diminishing} partial effects near the extremes.
  \end{keybox}

\end{frame}

\begin{frame}[shrink=8]{Where Do Logit and Probit Come From? A Latent Variable}

  \begin{definitionbox}{Latent variable model}
    Let $y^{*}$ be an unobserved (latent) ``propensity'':
    \[
      y^{*} = \beta_0 + \xbeta + e, \qquad y = \mathbf{1}[\,y^{*} > 0\,] .
    \]
    The error $e$ is \textbf{independent} of $\mathbf{x}$ and symmetric about zero.
  \end{definitionbox}

  \smallskip

  Then, using symmetry $1 - G(-z) = G(z)$,
  \[
    P(y=1\mid\mathbf{x}) = P(y^{*}>0\mid\mathbf{x})
      = P\!\big(e > -(\beta_0+\xbeta)\big) = G(\beta_0+\xbeta).
  \]

  \begin{keybox}
    \small
    \begin{itemize}
      \item $e \sim$ \textbf{logistic} $\Rightarrow$ logit; \quad
            $e \sim$ \textbf{normal} $\Rightarrow$ probit.
      \item The latent scale has \emph{no natural units}, so raw $\beta_j$ are not
        directly meaningful --- only their \textbf{sign} and
        \textbf{response-probability} effects.
    \end{itemize}
  \end{keybox}

\end{frame}

\begin{frame}[shrink=10]{Partial Effects: Why a Probit Coefficient Is Not a Marginal Effect}

  For a (roughly) continuous $x_j$, the chain rule gives
  \[
    \frac{\partial P(y=1\mid\mathbf{x})}{\partial x_j}
      = g(\beta_0 + \xbeta)\,\beta_j, \qquad g(z) \equiv \frac{dG}{dz}(z) > 0 .
  \]

  \begin{highlightbox}
    The effect of $x_j$ is $\beta_j$ \textbf{scaled} by the density $g(\beta_0+\xbeta)$,
    which \emph{depends on the values of all regressors}. So the partial effect
    \textbf{varies across observations}.
  \end{highlightbox}

  \smallskip

  \begin{columns}[t]
    \column{0.52\textwidth}
    For a \textbf{discrete} change in $x_1$ from $0$ to $1$:
    \[
      G(\beta_0+\beta_1+\beta_2 x_2 + \cdots) - G(\beta_0+\beta_2 x_2 + \cdots).
    \]
    \column{0.46\textwidth}
    \begin{keybox}
      \small
      \textbf{Anti-pattern:} ``the probit coefficient \emph{is} the marginal
      effect.'' \textcolor{negative}{Wrong} --- it is $\beta_j$ before scaling by
      $g(\cdot)$.
    \end{keybox}
  \end{columns}

  \smallskip

  \textbf{Relative} effects are clean, though:
  $\dfrac{\partial P/\partial x_j}{\partial P/\partial x_h} = \dfrac{\beta_j}{\beta_h}$
  --- the scale factor cancels.

\end{frame}

\begin{frame}[shrink=12]{Maximum Likelihood Estimation}

  OLS does not apply --- $G$ is nonlinear. We use \textbf{MLE}. With $y\mid\mathbf{x}$
  Bernoulli, the density for observation $i$ is
  \[
    f(y_i\mid\mathbf{x}_i) = \big[G(\xib)\big]^{y_i}\big[1-G(\xib)\big]^{1-y_i},
    \quad y_i \in \{0,1\}.
  \]

  \begin{definitionbox}{Log-likelihood}
    \[
      \ell_i(\boldsymbol{\beta}) = y_i \log\!\big[G(\xib)\big]
        + (1-y_i)\log\!\big[1 - G(\xib)\big],
      \qquad
      \mathcal{L}(\boldsymbol{\beta}) = \sum_{i=1}^{n} \ell_i(\boldsymbol{\beta}).
    \]
  \end{definitionbox}

  \smallskip

  \begin{keybox}
    \begin{itemize}
      \item $\hat{\boldsymbol\beta}$ maximizes $\mathcal{L}$; no closed form ---
        solved numerically.
      \item Under general conditions MLE is \textbf{consistent, asymptotically
        normal, and efficient}.
      \item Each $\hat\beta_j$ comes with an asymptotic SE $\Rightarrow$ form
        $t = \hat\beta_j / \se(\hat\beta_j)$ exactly as before.
    \end{itemize}
  \end{keybox}

\end{frame}

\begin{frame}[shrink=10]{First-Order Conditions: The Score Equations}

  $\hat{\boldsymbol\beta}$ solves $\partial\mathcal{L}/\partial\boldsymbol\beta=\mathbf{0}$.
  Differentiating $\ell_i = y_i\log G + (1-y_i)\log(1-G)$ with $g\equiv G'$:

  \begin{definitionbox}{Score (likelihood) equations}
    \[
      \sum_{i=1}^{n}
        \frac{\big[\,y_i - G(\xib)\,\big]\,g(\xib)}
             {G(\xib)\big[1-G(\xib)\big]}\;\mathbf{x}_i = \mathbf{0}.
    \]
  \end{definitionbox}

  \smallskip

  \begin{highlightbox}
    \textbf{Logit's clean special case.} For $G=\Lambda$ we have
    $g = \Lambda(1-\Lambda) = G(1-G)$, so the weight collapses to $1$:
    \[
      \sum_{i=1}^{n}\big[\,y_i - \Lambda(\xib)\,\big]\,\mathbf{x}_i = \mathbf{0}.
    \]
  \end{highlightbox}

  \smallskip

  \begin{keybox}
    \small
    The logit score is exactly ``residuals orthogonal to the regressors'' --- the
    MLE analog of the OLS normal equations. With an intercept in $\mathbf{x}_i$,
    the fitted probabilities average to the sample frequency:
    $n^{-1}\sum_i \Lambda(\xib) = \bar y$. No closed form exists; solve numerically
    (Newton--Raphson).
  \end{keybox}

\end{frame}

\begin{frame}[shrink=9]{Testing and Goodness-of-Fit}

  \begin{columns}[t]
    \column{0.5\textwidth}
    \begin{definitionbox}{Likelihood ratio test}
      For $q$ exclusion restrictions,
      \[
        LR = 2(\mathcal{L}_{ur} - \mathcal{L}_{r}) \;\overset{a}{\sim}\; \chi^2_q .
      \]
      The MLE analog of the $F$ test: dropping variables can only lower
      $\mathcal{L}$.
    \end{definitionbox}

    \column{0.5\textwidth}
    \begin{definitionbox}{McFadden's pseudo-$R^2$}
      \[
        \tilde\rho^2 = 1 - \frac{\mathcal{L}_{ur}}{\mathcal{L}_{0}},
      \]
      where $\mathcal{L}_0$ uses an intercept only. Lies in $[0,1)$ but is
      \textbf{not} a usual $R^2$.
    \end{definitionbox}
  \end{columns}

  \smallskip

  \begin{keybox}
    \textbf{Percent correctly predicted:} $\tilde y_i = 1$ if
    $G(\hat\beta_0 + \xib) \ge 0.5$. Useful, but \emph{misleading} when one
    outcome is rare --- report the percent correct \emph{within each} outcome
    ($\hat q_0,\hat q_1$), not just overall.

    \smallskip
    \textbf{Anti-pattern:} ``pseudo-$R^2 = R^2$.'' Not comparable; $0.2$--$0.4$ is
    already a respectable McFadden value.
  \end{keybox}

\end{frame}

\begin{frame}[shrink=10]{MROZ: Labor Force Participation --- LPM vs.\ Logit vs.\ Probit}

  \begin{examplebox}{Table 17.1 --- Dependent variable: $inlf$ (in labor force)}
    \centering\footnotesize
    \begin{tabular}{l c c c}
      \toprule
      Variable & LPM (OLS) & Logit (MLE) & Probit (MLE) \\
      \midrule
      $nwifeinc$ & $-0.0034$ $(0.0015)$ & $-0.021$ $(0.008)$ & $-0.012$ $(0.005)$ \\
      $educ$     & $0.038$ $(0.007)$    & $0.221$ $(0.043)$  & $0.131$ $(0.025)$  \\
      $exper$    & $0.039$ $(0.006)$    & $0.206$ $(0.032)$  & $0.123$ $(0.019)$  \\
      $exper^2$  & $-0.00060$ $(0.00019)$ & $-0.0032$ $(0.0010)$ & $-0.0019$ $(0.0006)$ \\
      $age$      & $-0.016$ $(0.002)$   & $-0.088$ $(0.015)$ & $-0.053$ $(0.008)$ \\
      $kidslt6$  & $-0.262$ $(0.032)$   & $-1.443$ $(0.204)$ & $-0.868$ $(0.119)$ \\
      $kidsge6$  & $0.013$ $(0.014)$    & $0.060$ $(0.075)$  & $0.036$ $(0.043)$ \\
      \midrule
      \% correct pred.\ / pseudo $R^2$ & $73.4$ / $0.264$ & $73.6$ / $0.220$ & $73.4$ / $0.221$ \\
      \bottomrule
    \end{tabular}
  \end{examplebox}

  \begin{keybox}
    \small
    Signs and significance agree across all three. \textbf{Magnitudes do not} ---
    coefficients live on different scales. To compare, we need \emph{partial
    effects}, not raw coefficients.
  \end{keybox}

\end{frame}

%% ============================================================
\section{Summarizing Effects: APE and PEA (\S17-1d)}
%% ============================================================

\begin{frame}[shrink=12]{Two Ways to Get One Number: PEA and APE}

  The partial effect $g(\beta_0+\xbeta)\beta_j$ varies by $i$. We summarize with
  a single \emph{scale factor}.

  \smallskip

  \begin{columns}[t]
    \column{0.49\textwidth}
    \begin{definitionbox}{Partial Effect at the Average (PEA)}
      Plug in the \textbf{average} regressors, then scale:
      \[
        \PEA_j = g(\hat\beta_0 + \xbar\hat{\boldsymbol\beta})\,\hat\beta_j .
      \]
      Problem: ``the average person'' may not exist (e.g.\ $\overline{female}=0.48$).
    \end{definitionbox}

    \column{0.49\textwidth}
    \begin{definitionbox}{Average Partial Effect (APE)}
      Scale \emph{each} person, then average:
      \[
        \APE_j = \Big[\tfrac{1}{n}\sum_{i=1}^{n} g(\hat\beta_0+\xib)\Big]\hat\beta_j .
      \]
      Also called the average marginal effect (AME). \textbf{Preferred} in practice.
    \end{definitionbox}
  \end{columns}

  \smallskip

  \begin{keybox}
    For \textbf{discrete} regressors, do not use the calculus approximation ---
    difference $G$ directly and average:
    \(
      n^{-1}\sum_i \big[G(\cdots + \hat\beta_k(c_k+1)) - G(\cdots + \hat\beta_k c_k)\big].
    \)
  \end{keybox}

\end{frame}

\begin{frame}[shrink=8]{MROZ: Average Partial Effects Reconcile the Models}

  \begin{examplebox}{Table 17.2 --- APEs for labor force participation}
    \centering\footnotesize
    \begin{tabular}{l c c c}
      \toprule
      Variable & LPM & Logit & Probit \\
      \midrule
      $nwifeinc$ & $-0.0034$ & $-0.0038$ & $-0.0036$ \\
      $educ$     & $0.038$   & $0.039$   & $0.039$ \\
      $exper$    & $0.027$   & $0.025$   & $0.026$ \\
      $age$      & $-0.016$  & $-0.016$  & $-0.016$ \\
      $kidslt6$  & $-0.262$  & $-0.258$  & $-0.261$ \\
      $kidsge6$  & $0.013$   & $0.011$   & $0.011$ \\
      \bottomrule
    \end{tabular}
  \end{examplebox}

  \begin{highlightbox}
    \small
    \textbf{The punchline.} Once converted to APEs, all three models tell nearly
    the \emph{same} story --- another year of schooling raises participation by
    about $3.9$ percentage points. The LPM was a decent approximation
    \emph{to the average effect} all along.

    \smallskip
    Rule of thumb: $\hat\beta^{\text{logit}}\approx 1.6\,\hat\beta^{\text{probit}}$;
    map to LPM by multiplying probit by $\approx0.4$, logit by $\approx0.25$.
  \end{highlightbox}

\end{frame}

\begin{frame}[shrink=10]{Diminishing Effects: Probit vs.\ LPM Across Education}

  \begin{columns}[t]
    \column{0.52\textwidth}
    \begin{center}
    \begin{tikzpicture}[scale=0.8]
      \draw[->,thick] (0,0) -- (6.2,0) node[below right] {\small years educ};
      \draw[->,thick] (0,0) -- (0,3.4) node[above] {\small $\widehat{P}(inlf=1)$};
      \foreach \y/\l in {0/0,1.5/0.5,3.0/1}
        \draw (-0.08,\y) -- (0.08,\y) node[left,font=\scriptsize] {\l};
      % LPM straight line: inlf = 0.102 + 0.038 educ, educ 0..20 -> p .102..0.862
      \draw[thick,SUFEGold] (0,0.31) -- (6,2.59);
      % probit S-curve Phi(-1.403 + 0.131 educ): crosses LPM ~ 11.5 yrs
      \draw[thick,SUFEBlue,domain=0:6,samples=80,smooth]
        plot (\x, {3.0/(1+exp(-(0.131*(\x*3.33) - 1.403)*2.2)) });
      \node[SUFEGold,font=\scriptsize] at (4.6,2.05) {LPM};
      \node[SUFEBlue,font=\scriptsize] at (4.4,2.95) {probit};
      \fill[black] (3.45,1.6) circle (1.5pt);
      \node[font=\scriptsize] at (3.45,1.30) {cross $\approx$ 11.5};
    \end{tikzpicture}
    \end{center}

    \column{0.46\textwidth}
    \begin{keybox}
      \small
      Recreates \textbf{Figure 17.2}.
      \begin{itemize}
        \item Low education: LPM predicts \emph{higher} participation than probit.
        \item Curves cross near $11.5$ years.
        \item High education: probit pulls ahead.
        \item LPM's constant slope is the cost of its simplicity.
      \end{itemize}
    \end{keybox}
  \end{columns}

  \smallskip

  \begin{highlightbox}
    \small
    \textbf{Answer to Q1.} A probit fitted value is $\Phi(\cdot)\in(0,1)$ by
    construction --- it can \emph{never} exceed $1$, unlike the LPM.
  \end{highlightbox}

\end{frame}

%% ============================================================
\section{Fractional Response Models (\S17-2)}
%% ============================================================

\begin{frame}[shrink=8]{Fractional Outcomes: Building the Model}

  A \textbf{fractional response} takes values in $[0,1]$, possibly including the
  endpoints: the fraction of employees in a pension plan, the share of votes won,
  the proportion of students who pass an exam.

  \smallskip

  \begin{keybox}
    The object of interest is the \textbf{conditional mean} $\E(y\mid\mathbf{x})$,
    which must itself lie in $[0,1]$. A linear model
    $\E(y\mid\mathbf{x})=\beta_0+\xbeta$ can produce means outside $[0,1]$ ---
    exactly the LPM defect, now for a \emph{mean} rather than a probability.
  \end{keybox}

  \smallskip

  \begin{definitionbox}{Fractional response model}
    Pass the index through a CDF $G$ with $0 < G(z) < 1$:
    \[
      \E(y\mid\mathbf{x}) = G(\beta_0 + \xbeta), \qquad y \in [0,1].
    \]
    \textbf{Fractional logit:}\;\;
    $\E(y\mid\mathbf{x}) = \Lambda(\beta_0+\xbeta)
       = \dfrac{\exp(\beta_0+\xbeta)}{1+\exp(\beta_0+\xbeta)}$.
    \\[3pt]
    \textbf{Fractional probit:}\;\;
    $\E(y\mid\mathbf{x}) = \Phi(\beta_0+\xbeta)$.
  \end{definitionbox}

  \smallskip

  \begin{keybox}
    \small
    The \emph{same} functional forms as binary logit/probit --- but here $G(\cdot)$
    is a conditional \textbf{mean}, not a probability. Binary response is the
    special case $y\in\{0,1\}$, where the mean \emph{equals} $P(y=1\mid\mathbf{x})$.
  \end{keybox}

\end{frame}

\begin{frame}[shrink=10]{Fractional Logit: Quasi-Maximum Likelihood Estimation}

  We cannot use a Bernoulli \emph{likelihood} --- $y$ is no longer in $\{0,1\}$.
  Yet Papke and Wooldridge \citep{PapkeWooldridge1996_fractional} show that we may
  maximize the \emph{same} Bernoulli quasi-log-likelihood:

  \begin{definitionbox}{Bernoulli quasi-log-likelihood}
    \[
      \ell_i(\boldsymbol\beta) = y_i\log\!\big[G(\xib)\big]
        + (1-y_i)\log\!\big[1-G(\xib)\big], \qquad y_i \in [0,1].
    \]
    Maximizing $\sum_i \ell_i(\boldsymbol\beta)$ gives the \textbf{fractional logit}
    ($G=\Lambda$) or \textbf{fractional probit} ($G=\Phi$) estimator.
  \end{definitionbox}

  \smallskip

  \begin{keybox}
    \small
    \begin{itemize}
      \item \textbf{Quasi}-MLE: the likelihood is knowingly \emph{wrong} ($y$ is not
        Bernoulli), hence ``QMLE'' --- yet $\hat{\boldsymbol\beta}$ is consistent.
      \item \textbf{Consistency} requires only that the mean
        $\E(y\mid\mathbf{x})=G(\beta_0+\xbeta)$ be correct; the rest of the
        distribution is left unspecified.
      \item \textbf{Robust SEs} are mandatory: the Bernoulli variance $G(1-G)$ is
        incorrect for $y\notin\{0,1\}$ (it overstates, so robust SEs are smaller).
      \item \textbf{Partial effects} as in binary:
        $\partial\E(y\mid\mathbf{x})/\partial x_j = g(\beta_0+\xbeta)\,\beta_j$, with
        $g(z)=\Lambda(z)[1-\Lambda(z)]$ for fractional logit.
    \end{itemize}
  \end{keybox}

\end{frame}

\begin{frame}{Example: 401(k) Participation Rates}

  \begin{center}
    {\small\textbf{\textcolor{SUFEBlue}{Table 17.3 --- $prate$ participation rate ($n=4075$)}}}

    \smallskip
    \footnotesize
    \begin{tabular}{l c c c}
      \toprule
      Variable & Linear OLS & Frac.\ Logit & Frac.\ Logit APE \\
      \midrule
      $mrate$  & $0.107$ $(0.006)$  & $1.158$ $(0.075)$  & $0.147$ $(0.010)$ \\
      $age$    & $0.0037$ $(0.0002)$ & $0.035$ $(0.003)$ & $0.0044$ $(0.0004)$ \\
      $ltotemp$& $-0.028$ $(0.002)$ & $-0.207$ $(0.014)$ & $-0.026$ $(0.002)$ \\
      $sole$   & $0.018$ $(0.006)$  & $0.166$ $(0.051)$  & $0.021$ $(0.006)$ \\
      \midrule
      $R^2$    & $0.176$ & $0.195$ & --- \\
      \bottomrule
    \end{tabular}
  \end{center}

  \smallskip

  \begin{keybox}
    \small
    The linear model produces $159$ fitted values \emph{above one}; fractional
    logit cannot. Its match-rate APE ($0.147$) exceeds the OLS slope ($0.107$),
    capturing a \emph{diminishing} effect a linear model needs an ad-hoc
    $mrate^{2}$ to mimic.
  \end{keybox}

\end{frame}

%% ============================================================
\section{Nonnegative Outcomes: Poisson Regression (\S17-3)}
%% ============================================================

\begin{frame}[shrink=12]{Counts and Nonnegative $y$: The Exponential Mean}

  When $y \in \{0,1,2,\ldots\}$ (arrests, patents) or any $y\ge 0$, a linear model
  can predict negatives. Use an \textbf{exponential conditional mean}:
  \[
    \E(y\mid\mathbf{x}) = \exp(\beta_0 + \xbeta) > 0 .
  \]

  \begin{highlightbox}
    Taking logs, $\log \E(y\mid\mathbf{x}) = \beta_0 + \xbeta$, so
    $100\,\beta_j$ is approximately the \textbf{percentage change} in
    $\E(y\mid\mathbf{x})$ for a one-unit rise in $x_j$ --- familiar
    semi-elasticity interpretation, but \emph{valid even when} $y=0$.
  \end{highlightbox}

  \smallskip

  \begin{keybox}
    \textbf{Why not just use $\log(1+y)$?} Because it is \emph{not} scale-invariant
    in $y$ and does not give clean percentage-change coefficients
    \citep{ChenRoth2024_logs}. The exponential mean sidesteps the ``log of zero''
    problem entirely.
  \end{keybox}

\end{frame}

\begin{frame}[shrink=10]{Poisson Regression and the QMLE Escape Hatch}

  \begin{definitionbox}{Poisson model}
    Assume $y\mid\mathbf{x}$ is Poisson with mean $\exp(\xbeta)$:
    \[
      P(y=h\mid\mathbf{x}) = \frac{\exp[-\exp(\xbeta)]\,[\exp(\xbeta)]^{h}}{h!},
      \quad h = 0,1,2,\ldots
    \]
    Maximize $\mathcal{L}(\boldsymbol\beta)=\sum_i\{y_i\xib - \exp(\xib)\}$.
  \end{definitionbox}

  \smallskip

  The Poisson assumption forces $\Var(y\mid\mathbf{x}) = \E(y\mid\mathbf{x})$ ---
  often violated (\textbf{overdispersion}, $\sigma^2>1$).

  \begin{keybox}
    \small
    \textbf{Poisson QMLE robustness.} $\hat{\boldsymbol\beta}$ is \emph{consistent
    and asymptotically normal whenever the mean $\exp(\xbeta)$ is correct} ---
    regardless of the true variance, even if $y$ is not a count. Just inflate SEs
    by $\hat\sigma$, with $\hat\sigma^2 = (n-k-1)^{-1}\sum_i \hat u_i^2/\hat y_i$.
    This is why Poisson beats negative binomial \citep{SantosSilvaTenreyro2006_log}.
  \end{keybox}

\end{frame}

\begin{frame}[shrink=8]{Example: Number of Arrests (\texttt{CRIME1})}

  \begin{examplebox}{Table 17.5 --- Dependent variable: $narr86$ (arrests in 1986)}
    \centering\footnotesize
    \begin{tabular}{l c c}
      \toprule
      Variable & Linear (OLS) & Exponential (Poisson QMLE) \\
      \midrule
      $pcnv$   & $-0.132$ $(0.040)$ & $-0.402$ $(0.085)$ \\
      $ptime86$& $-0.041$ $(0.009)$ & $-0.099$ $(0.021)$ \\
      $inc86$  & $-0.0015$ $(0.0003)$ & $-0.0081$ $(0.0010)$ \\
      $black$  & $0.327$ $(0.045)$  & $0.661$ $(0.074)$ \\
      \midrule
      Log-likelihood & --- & $-2{,}248.76$ \\
      $\hat\sigma$ & --- & $1.232$ \\
      \bottomrule
    \end{tabular}
  \end{examplebox}

  \begin{keybox}
    \small
    $narr86$ is $0$ for $1{,}970$ of $2{,}725$ men --- a count with a huge zero
    pileup, ideal for Poisson. A higher conviction probability
    ($\Delta pcnv = 0.1$) lowers expected arrests by about $4\%$; a black man's
    expected arrests are $100[\exp(0.661)-1]\approx 94\%$ higher, ceteris paribus.
    Since $\hat\sigma = 1.232 > 1$ (overdispersion), inflate SEs by $\approx 23\%$.
  \end{keybox}

\end{frame}

%% ============================================================
\section{Corner Solutions: The Tobit Model (\S17-4)}
%% ============================================================

\begin{frame}{Corner Solutions: A Pileup of Zeros}

  \begin{columns}[c]
    \column{0.55\textwidth}
    Hours worked, charitable gifts, insurance: a \textbf{nontrivial fraction
    choose exactly zero}, then a continuous spread of positives. OLS predicts
    negatives, cannot match the spike at zero, and is heteroskedastic.

    \column{0.43\textwidth}
    \centering
    \begin{tikzpicture}[scale=0.7]
      \draw[->,thick] (0,0) -- (4.6,0) node[below right] {\scriptsize $y$};
      \draw[->,thick] (0,0) -- (0,2.9) node[above] {\scriptsize density};
      \draw[very thick,negative] (0.12,0) -- (0.12,2.4);
      \node[negative,font=\scriptsize,right] at (0.18,2.4) {mass at $0$};
      \draw[thick,SUFEBlue,domain=0.6:4.3,samples=80,smooth]
        plot (\x, {2.0*exp(-((\x-2.1)^2)/1.4)});
      \node[SUFEBlue,font=\scriptsize] at (3.1,1.8) {positive $y$};
    \end{tikzpicture}
  \end{columns}

  \smallskip

  \begin{definitionbox}{Tobit model \citep{Tobin1958_estimation}}
    A latent linear model, observed only when positive:
    \[
      y^{*} = \beta_0 + \xbeta + u, \quad u\mid\mathbf{x}\sim \Normal(0,\sigma^2),
      \qquad y = \max(0,\, y^{*}).
    \]
  \end{definitionbox}

\end{frame}

\begin{frame}[shrink=8]{Tobit: Likelihood and the Role of $\sigma$}

  Two pieces: a \textbf{probability} of the corner and a \textbf{density} for positives.
  \[
    P(y=0\mid\mathbf{x}) = 1 - \Phi(\xbeta/\sigma), \qquad
    y>0:\; \tfrac{1}{\sigma}\phi\!\big((y-\xbeta)/\sigma\big).
  \]

  \begin{definitionbox}{Tobit log-likelihood}
    \[
      \ell_i(\boldsymbol\beta,\sigma)
       = \mathbf{1}[y_i=0]\,\log\!\big[1-\Phi(\xib/\sigma)\big]
       + \mathbf{1}[y_i>0]\,\log\!\Big\{\tfrac{1}{\sigma}\phi\big((y_i-\xib)/\sigma\big)\Big\}.
    \]
  \end{definitionbox}

  \smallskip

  \begin{keybox}
    Estimate $\boldsymbol\beta$ \emph{and} $\sigma$ jointly by MLE.
    \textbf{$\sigma$ is not ``ancillary''} --- it enters the partial effects
    directly, so it matters for the economics, not just for inference.
  \end{keybox}

\end{frame}

\begin{frame}[shrink=8]{The Inverse Mills Ratio: A Truncated-Normal Lemma}

  Every Tobit (and selection) expectation rests on one fact about the normal.

  \begin{definitionbox}{Truncated-normal lemma}
    If $z\sim\Normal(0,1)$, then for any constant $c$,
    \[
      \E\big[z \mid z > c\big] = \frac{\phi(c)}{1-\Phi(c)},
      \qquad
      \E\big[z \mid z > -c\big] = \frac{\phi(c)}{\Phi(c)} \;\equiv\; \lambda(c),
    \]
    where the second form uses $\phi(-c)=\phi(c)$ and $1-\Phi(-c)=\Phi(c)$.
  \end{definitionbox}

  \smallskip

  \begin{keybox}
    \textbf{Inverse Mills ratio} $\lambda(c)=\phi(c)/\Phi(c)$ --- the ``selection hazard'':
    \begin{itemize}
      \item $\lambda(c) > 0$ for all $c$, and is \textbf{strictly decreasing}.
      \item $\lambda(c)\to 0$ as $c\to+\infty$; $\lambda(c)\approx -c$ as
        $c\to-\infty$, so $c+\lambda(c)>0$ always.
      \item Derivative: $\lambda'(c) = -\lambda(c)\big[c+\lambda(c)\big]\in(-1,0)$.
    \end{itemize}
  \end{keybox}

  \smallskip

  \begin{highlightbox}
    \small
    This one function does the work in \emph{both} the Tobit conditional mean
    $\E(y\mid y>0,\mathbf{x})$ and the Heckman selection correction. Watch for it twice.
  \end{highlightbox}

\end{frame}

\begin{frame}[shrink=8]{Tobit Partial Effects via the Inverse Mills Ratio}

  The $\beta_j$ measure effects on the \emph{latent} $y^{*}$ --- rarely the goal;
  we want effects on the \textbf{observed} $y$. Apply the lemma with $z=u/\sigma$,
  since $y>0 \Leftrightarrow u/\sigma > -\xbeta/\sigma$:
  \[
    \E(y\mid y>0,\mathbf{x})
      = \xbeta + \sigma\,\E\!\Big[\tfrac{u}{\sigma}\,\Big|\,
        \tfrac{u}{\sigma} > -\tfrac{\xbeta}{\sigma}\Big]
      = \xbeta + \sigma\,\lambda(\xbeta/\sigma).
  \]
  Combine with $P(y>0\mid\mathbf{x})=\Phi(\xbeta/\sigma)$; the $\lambda$ term cancels:
  \[
    \E(y\mid\mathbf{x}) = \Phi(\xbeta/\sigma)\,\E(y\mid y>0,\mathbf{x})
      = \Phi(\xbeta/\sigma)\,\xbeta + \sigma\,\phi(\xbeta/\sigma).
  \]

  \begin{keybox}
    \small
    Hence the partial effect on the \textbf{unconditional} mean is remarkably clean:
    \[
      \frac{\partial\E(y\mid\mathbf{x})}{\partial x_j} = \beta_j\,\Phi(\xbeta/\sigma).
    \]
    The factor $\Phi(\xbeta/\sigma)\in(0,1)$ scales the latent coefficient down to
    the observed-mean effect --- the APE scale factor for Tobit.
  \end{keybox}

\end{frame}

\begin{frame}[shrink=8]{MROZ: Annual Hours Worked --- OLS vs.\ Tobit}

  \begin{examplebox}{Table 17.6/17.7 --- Dependent variable: $hours$ ($325$ zeros)}
    \centering\footnotesize
    \begin{tabular}{l c c c c}
      \toprule
       & \multicolumn{2}{c}{Coefficients} & \multicolumn{2}{c}{APE} \\
      Variable & OLS & Tobit & OLS & Tobit \\
      \midrule
      $nwifeinc$ & $-3.45$ $(2.24)$ & $-8.81$ $(4.46)$ & $-3.45$ & $-5.19$ \\
      $educ$     & $28.76$ $(13.04)$ & $80.65$ $(21.58)$ & $28.76$ & $47.47$ \\
      $age$      & $-30.51$ $(4.24)$ & $-54.41$ $(7.42)$ & $-30.51$ & $-32.03$ \\
      $kidslt6$  & $-442.1$ $(57.5)$ & $-894.0$ $(111.9)$ & $-442.1$ & $-526.3$ \\
      \midrule
      $\hat\sigma$ & $750.18$ & $1{,}122.02$ & & \\
      \bottomrule
    \end{tabular}
  \end{examplebox}

  \begin{keybox}
    \small
    \textbf{Do not compare raw coefficients.} The Tobit $educ$ coefficient
    ($80.65$) dwarfs OLS ($28.76$). But the APE scale factor
    $n^{-1}\sum_i\Phi(\xib/\hat\sigma)\approx 0.589$ gives a Tobit APE of
    $0.589\times 80.65 \approx 47.5$ hours --- larger than OLS, but far less
    dramatic than the raw coefficient suggests.
  \end{keybox}

\end{frame}

\begin{frame}[shrink=8]{When Is Tobit the Wrong Model?}

  \begin{highlightbox}
    Tobit ties $P(y>0\mid\mathbf{x})$ and $\E(y\mid y>0,\mathbf{x})$ to the
    \emph{same} parameters $\boldsymbol\beta/\sigma$. A regressor cannot raise the
    chance of a positive outcome while lowering its size.
  \end{highlightbox}

  \smallskip

  \begin{keybox}
    \textbf{Specification check.} Run a probit of $w=\mathbf{1}[y>0]$ on
    $\mathbf{x}$; its coefficients should be ``close'' to $\hat\beta_j/\hat\sigma$
    from Tobit. Sign flips or large magnitude gaps signal that Tobit is too
    restrictive --- consider a two-part (hurdle) model instead.

    \smallskip
    \textbf{Answer to Q2.} In \texttt{MROZ}, ``hours'' is a \textbf{corner
    solution}, not censoring --- women genuinely \emph{choose} zero. Tobit is
    appropriate, and $\sigma$ is economically meaningful (not a nuisance).
  \end{keybox}

\end{frame}

%% ============================================================
\section{Censored and Truncated Regression (\S17-5)}
%% ============================================================

\begin{frame}[shrink=10]{Censoring vs.\ Truncation: A Data Problem, Not a Behavior}

  \begin{columns}[t]
    \column{0.49\textwidth}
    \begin{definitionbox}{Censored regression}
      We \textbf{always observe $\mathbf{x}_i$} for a random draw, but $y_i$ is
      only recorded up to a threshold $c_i$:
      \[
        y_i = \beta_0 + \xib + u_i, \quad w_i = \min(y_i, c_i).
      \]
      Example: top-coded wealth (``$> \$500\text{k}$''), duration data.
    \end{definitionbox}

    \column{0.49\textwidth}
    \begin{definitionbox}{Truncated regression}
      A \textbf{sub-population is excluded entirely} --- no $\mathbf{x}_i$, no
      $y_i$ --- based on $y$:
      \[
        \text{observe }(\mathbf{x}_i,y_i)\text{ only if } y_i \le c_i.
      \]
      Example: a study that only enrolls families below 1.5$\times$ the poverty line.
    \end{definitionbox}
  \end{columns}

  \smallskip

  \begin{keybox}
    \textbf{Key difference.} Censoring keeps the unit (we know $y>c$);
    truncation \emph{drops} it. Censored regression exploits the censored units;
    truncated regression cannot.
  \end{keybox}

\end{frame}

\begin{frame}[shrink=8]{Censored Normal Regression: Likelihood}

  Under $y_i = \beta_0 + \xib + u_i$, $u_i\mid\mathbf{x}_i,c_i \sim \Normal(0,\sigma^2)$
  and $w_i=\min(y_i,c_i)$:

  \begin{definitionbox}{Density of the observed $w_i$}
    \[
      f(w\mid\mathbf{x}_i,c_i) =
      \begin{cases}
        1 - \Phi\!\big((c_i - \xib)/\sigma\big), & w = c_i \quad(\text{censored}),\\[4pt]
        \tfrac{1}{\sigma}\phi\!\big((w - \xib)/\sigma\big), & w < c_i \quad(\text{uncensored}).
      \end{cases}
    \]
  \end{definitionbox}

  \smallskip

  \begin{keybox}
    \begin{itemize}
      \item Maximize $\sum_i \log f(\cdot)$ to get $\hat{\boldsymbol\beta},\hat\sigma$.
      \item \textbf{Crucial:} here the $\beta_j$ are interpreted \emph{just as in a
        linear regression} --- no inverse-Mills gymnastics. The expectations of
        interest are linear in $\boldsymbol\beta$.
      \item Naively running OLS on \emph{uncensored} obs only $\Rightarrow$
        inconsistent (same algebra as Tobit, different cause).
    \end{itemize}
  \end{keybox}

\end{frame}

\begin{frame}[shrink=10]{Example: Duration of Recidivism (\texttt{RECID})}

  \begin{columns}[t]
    \column{0.55\textwidth}
    \begin{examplebox}{Table 17.8 --- DV $\log(durat)$ --- $62\%$ censored}
      \centering\footnotesize
      \begin{tabular}{l c}
        \toprule
        Variable & Coef.\ (SE) \\
        \midrule
        $workprg$ & $-0.063$ $(0.120)$ \\
        $priors$  & $-0.137$ $(0.021)$ \\
        $tserved$ & $-0.019$ $(0.003)$ \\
        $felon$   & $0.444$ $(0.145)$ \\
        $alcohol$ & $-0.635$ $(0.144)$ \\
        $married$ & $0.341$ $(0.140)$ \\
        \midrule
        $\hat\sigma$ & $1.810$ \\
        \bottomrule
      \end{tabular}
    \end{examplebox}

    \column{0.43\textwidth}
    \begin{keybox}
      \small
      $durat$ = months to re-arrest; $893$ of $1{,}445$ never re-arrested
      $\Rightarrow$ \textbf{censored}.
      \begin{itemize}
        \item One more prior $\Rightarrow$ $\approx 14\%$ shorter duration.
        \item $alcohol$ abuse sharply shortens duration.
        \item $workprg$ effect small, insignificant.
      \end{itemize}
    \end{keybox}
  \end{columns}

  \smallskip

  \begin{highlightbox}
    \small
    Ignoring censoring (OLS treating censored as actual) shrinks estimates toward
    zero: $priors$ goes from $-0.137$ to $-0.059$. Censored regression recovers the
    correct magnitudes.
  \end{highlightbox}

\end{frame}

\begin{frame}{Truncation Biases OLS Toward Zero}

  \begin{columns}[c]
    \column{0.52\textwidth}
    \centering
    \begin{tikzpicture}[scale=0.72]
      \draw[->,thick] (0,0) -- (5.6,0) node[below right] {\scriptsize education};
      \draw[->,thick] (0,0) -- (0,3.8) node[above] {\scriptsize income};
      % truncation line at income = 50
      \draw[dashed,gray] (0,2.3) -- (5.4,2.3) node[right,font=\scriptsize] {$c=50$};
      % observed points (below line) - solid
      \fill[SUFEBlue] (0.8,0.9) circle (2pt); \fill[SUFEBlue] (1.4,1.3) circle (2pt);
      \fill[SUFEBlue] (2.0,1.0) circle (2pt); \fill[SUFEBlue] (2.6,1.8) circle (2pt);
      \fill[SUFEBlue] (3.2,1.5) circle (2pt); \fill[SUFEBlue] (3.8,2.1) circle (2pt);
      % dropped points (above line) - hollow
      \draw[SUFEGold,thick] (3.0,2.7) circle (2pt); \draw[SUFEGold,thick] (3.6,3.0) circle (2pt);
      \draw[SUFEGold,thick] (4.2,2.6) circle (2pt); \draw[SUFEGold,thick] (4.6,3.3) circle (2pt);
      \draw[SUFEGold,thick] (5.0,2.9) circle (2pt);
      % true regression line (steep)
      \draw[thick,SUFEGold] (0.4,0.5) -- (5.2,3.4);
      \node[SUFEGold,font=\scriptsize] at (4.9,3.6) {true};
      % truncated (flattened) line
      \draw[thick,SUFEBlue] (0.4,0.8) -- (5.2,2.15);
      \node[SUFEBlue,font=\scriptsize] at (5.0,1.9) {truncated};
    \end{tikzpicture}

    \column{0.45\textwidth}
    \begin{keybox}
      \small
      Recreates \textbf{Figure 17.4}. Dropping incomes above $\$50\text{k}$
      (hollow points) lops off the upper-right tail; the fitted line through the
      survivors \textbf{flattens} --- biased toward zero.
    \end{keybox}

    \smallskip

    \begin{highlightbox}
      \small
      The \textbf{truncated normal MLE} restores consistency by re-normalizing the
      density over the observed region \citep{HausmanWise1977_social}.
    \end{highlightbox}
  \end{columns}

\end{frame}

%% ============================================================
\section{Sample Selection Corrections (\S17-6)}
%% ============================================================

\begin{frame}[shrink=10]{When Is OLS on a Selected Sample Still OK?}

  Let $s_i=1$ if observation $i$ is used, $s_i=0$ otherwise. We run OLS on the
  $s_i=1$ sub-sample.

  \smallskip

  \begin{columns}[t]
    \column{0.49\textwidth}
    \begin{highlightbox}
      \textbf{\textcolor{positive}{Exogenous selection --- OLS is fine}}
      \begin{itemize}
        \item $s$ depends only on $\mathbf{x}$ (or on factors independent of $u$)
        \item Random sub-sampling
        \item $\Rightarrow$ $\hat{\boldsymbol\beta}$ still consistent / unbiased
      \end{itemize}
    \end{highlightbox}

    \column{0.49\textwidth}
    \begin{highlightbox}
      \textbf{\textcolor{negative}{Endogenous selection --- OLS biased}}
      \begin{itemize}
        \item $s$ depends on $y$ (truncation on the outcome)
        \item $s$ depends on $u$ (unobservables)
        \item $\Rightarrow$ selection bias
      \end{itemize}
    \end{highlightbox}
  \end{columns}

  \smallskip

  \begin{keybox}
    \textbf{Rule.} Selecting on the explanatory variables is harmless. Selecting
    on the dependent variable (or on the error) is poison. The danger case is
    \textbf{incidental truncation}.
  \end{keybox}

\end{frame}

\begin{frame}[shrink=8]{Incidental Truncation: The Wage-Offer Problem}

  We want the \textbf{wage offer} equation for \emph{all} women, but we observe a
  wage only for those who \emph{choose to work}.

  \smallskip

  \begin{definitionbox}{Selection model \citep{Heckman1979_selection}}
    \begin{align*}
      \text{outcome:}\quad & y = \xbeta + u, \qquad \E(u\mid\mathbf{x},\mathbf{z}) = 0, \\
      \text{selection:}\quad & s = \mathbf{1}[\mathbf{z}\boldsymbol\gamma + v \ge 0], \quad v\sim\Normal(0,1).
    \end{align*}
    We observe $y$ only when $s=1$. Truncation of $y$ is \emph{incidental} --- it
    depends on another variable (working), not on $y$ itself.
  \end{definitionbox}

  \smallskip

  \begin{keybox}
    If $\Corr(u,v) = \rho \ne 0$, the working sample is \textbf{not}
    representative: women who work have systematically different $u$. OLS on
    workers is then inconsistent for $\boldsymbol\beta$.
  \end{keybox}

\end{frame}

\begin{frame}[shrink=8]{The Selection Bias as an Omitted Variable}

  Take the expectation of $y=\xbeta+u$ given $(\mathbf{z},v)$. With
  $\mathbf{x}\subseteq\mathbf{z}$, $(u,v)$ independent of $\mathbf{z}$, and $(u,v)$
  jointly normal so that $\E(u\mid v)=\rho v$:
  \[
    \E(y\mid\mathbf{z},v) = \xbeta + \E(u\mid v) = \xbeta + \rho v.
  \]
  Now condition on selection, $s=1 \Leftrightarrow v > -\mathbf{z}\boldsymbol\gamma$.
  By the \textbf{truncated-normal lemma}, $\E(v\mid v>-\mathbf{z}\boldsymbol\gamma)
  = \lambda(\mathbf{z}\boldsymbol\gamma)$:
  \[
    \E(y\mid\mathbf{z}, s=1) = \xbeta + \rho\,\lambda(\mathbf{z}\boldsymbol\gamma),
    \qquad \lambda(c) = \frac{\phi(c)}{\Phi(c)} .
  \]

  \begin{highlightbox}
    \small
    OLS of $y$ on $\mathbf{x}$ using workers only \textbf{omits} the regressor
    $\rho\,\lambda(\mathbf{z}\boldsymbol\gamma)$. Since $\lambda(\cdot)$ correlates
    with $\mathbf{x}$, this is classic \textbf{omitted-variable bias} --- vanishing
    only when $\rho=0$.
  \end{highlightbox}

  \begin{keybox}
    \small
    \textbf{The fix writes itself:} \emph{construct} the missing regressor
    $\lambda(\mathbf{z}\boldsymbol\gamma)$ and add it back. OLS on the selected
    sample is then consistent --- precisely Heckman's two-step idea.
  \end{keybox}

\end{frame}

\begin{frame}[shrink=10]{The Heckit Procedure}

  \begin{definitionbox}{Heckman two-step (Heckit)}
    \begin{enumerate}
      \item \textbf{Selection equation.} Using \emph{all} $n$ observations, estimate a
        probit of $s_i$ on $\mathbf{z}_i$. Form the estimated inverse Mills ratio
        \[
          \hat\lambda_i = \lambda(\mathbf{z}_i\hat{\boldsymbol\gamma})
            = \phi(\mathbf{z}_i\hat{\boldsymbol\gamma})/\Phi(\mathbf{z}_i\hat{\boldsymbol\gamma}).
        \]
      \item \textbf{Outcome equation.} Using only the selected sample ($s_i=1$), run
        \[
          y_i \;\text{ on }\; \mathbf{x}_i,\; \hat\lambda_i .
        \]
    \end{enumerate}
  \end{definitionbox}

  \smallskip

  \begin{keybox}
    \begin{itemize}
      \item The $t$ statistic on $\hat\lambda_i$ is a \textbf{direct test for
        selection bias} ($H_0:\rho=0$).
      \item \textbf{Exclusion restriction:} $\mathbf{z}$ should contain at least one
        variable \emph{not} in $\mathbf{x}$ --- something that shifts the decision to
        work but not the wage offer. Without it, multicollinearity inflates the SEs.
    \end{itemize}
  \end{keybox}

\end{frame}

\begin{frame}[shrink=12]{MROZ: Wage Offer Equation --- OLS vs.\ Heckit}

  \begin{examplebox}{Table 17.9 --- DV $\log(wage)$ --- $428$ workers}
    \centering\footnotesize
    \begin{tabular}{l c c}
      \toprule
      Variable & OLS & Heckit \\
      \midrule
      $educ$    & $0.108$ $(0.014)$ & $0.109$ $(0.016)$ \\
      $exper$   & $0.042$ $(0.012)$ & $0.044$ $(0.016)$ \\
      $exper^2$ & $-0.00081$ $(0.00039)$ & $-0.00086$ $(0.00044)$ \\
      $\hat\lambda$ & --- & $0.032$ $(0.134)$ \\
      \midrule
      $n$ & $428$ & $428$ \\
      \bottomrule
    \end{tabular}
  \end{examplebox}

  \begin{highlightbox}
    \small
    Excluded from $\mathbf{x}$ but in $\mathbf{z}$: $nwifeinc$, $age$, $kidslt6$,
    $kidsge6$ (drive participation, assumed not the wage offer).

    \smallskip
    \textbf{Answer to Q3.} Here $t_{\hat\lambda}=0.032/0.134 \approx 0.24$ --- we
    \emph{fail to reject} $\rho=0$. No evidence of selection bias; OLS and Heckit
    returns to education differ by one-tenth of a point. Testing first saved an
    unnecessary correction.
  \end{highlightbox}

\end{frame}

%% ============================================================
\section{Synthesis}
%% ============================================================

\begin{frame}{Decision Guide: Which Model for Which $y$?}

  \begin{center}
  \small
  \begin{tabular}{l l l}
    \toprule
    Nature of $y$ & Model & Key interpretation tool \\
    \midrule
    Binary $\{0,1\}$        & Logit / Probit (MLE)      & APE / PEA \\
    Fractional $[0,1]$      & Fractional logit (QMLE)   & APE, robust SE \\
    Count / nonneg.\ $y\ge0$& Poisson (QMLE)            & $100\beta_j\approx \%\Delta\E(y)$ \\
    Corner solution $y\ge0$ & Tobit (MLE)               & APE $=\Phi(\xbeta/\sigma)\beta_j$ \\
    Top-coded / duration    & Censored normal reg.      & $\beta_j$ as in linear model \\
    Sub-population excluded & Truncated normal reg.     & $\beta_j$ as in linear model \\
    Selection on outcome    & Heckit (two-step)         & inverse Mills + $t$-test on $\rho$ \\
    \bottomrule
  \end{tabular}
  \end{center}

  \smallskip

  \begin{keybox}
    \textbf{First question to ask:} is $y$ \emph{genuinely} limited in the
    population (behavior $\to$ logit/probit/Poisson/Tobit), or is it fine but
    \emph{badly observed} (data $\to$ censored/truncated/Heckit)?
  \end{keybox}

\end{frame}

\begin{frame}{Recurring Themes Across All Models}

  \begin{highlightbox}
    \textbf{1. Coefficients $\ne$ effects.} In every nonlinear model the raw
    $\beta_j$ must be transformed (APE, percentage change, Mills adjustment) before
    it answers an economic question.
  \end{highlightbox}

  \smallskip

  \begin{highlightbox}
    \textbf{2. MLE and QMLE do the heavy lifting.} Logit, probit, Tobit, censored,
    and truncated models are all MLE; fractional and Poisson are \emph{quasi}-MLE
    --- consistent for the mean even when the full distribution is wrong, provided
    SEs are made robust.
  \end{highlightbox}

  \smallskip

  \begin{highlightbox}
    \textbf{3. Selection is omitted-variable bias.} The inverse Mills ratio is the
    missing regressor; estimate it, add it back, and test whether it mattered.
  \end{highlightbox}

\end{frame}

\begin{frame}{Summary}

  \begin{keybox}
    \textbf{Binary \& fractional (\S17-1, \S17-2).} Logit/probit squash the index
    into $(0,1)$; estimate by MLE/QMLE; summarize with APEs. The LPM approximates
    the average effect well.

    \smallskip
    \textbf{Nonnegative (\S17-3).} Exponential mean + Poisson QMLE handles counts
    and zeros, gives percentage-change coefficients, and is robust to overdispersion.

    \smallskip
    \textbf{Corner solutions (\S17-4).} Tobit models the zero pileup; partial
    effects scale the latent coefficients by $\Phi(\xbeta/\sigma)$.

    \smallskip
    \textbf{Data problems (\S17-5, \S17-6).} Censored/truncated regression fix
    observability defects; Heckit corrects incidental-truncation selection bias and
    provides a built-in test for it.
  \end{keybox}

\end{frame}

\begin{frame}[shrink=12]{Bridge: From LDVs to Causal Inference}

  \begin{columns}[t]
    \column{0.48\textwidth}
    \begin{highlightbox}
      \textbf{Lecture 17 (today)} \\
      Modeling \emph{limited} outcomes:
      \begin{itemize}
        \item Respect the range of $y$
        \item Recover partial effects, not just coefficients
        \item Correct for non-random observation
      \end{itemize}
    \end{highlightbox}

    \column{0.48\textwidth}
    \begin{highlightbox}
      \textbf{Next: Causal Inference (Ch.\ 19)} \\
      Designs for credible effects:
      \begin{itemize}
        \item Treatment effects with binary/limited outcomes
        \item RD, DiD extensions, matching
        \item Potential outcomes meet nonlinear models
      \end{itemize}
    \end{highlightbox}
  \end{columns}

  \smallskip

  \begin{keybox}
    The logit/probit and Tobit machinery returns in causal settings: average
    treatment effects with nonlinear models are just APEs of a treatment dummy.
    The threads from today run straight into the next lecture.
  \end{keybox}

\end{frame}

\begin{frame}[allowframebreaks]{References}
  \footnotesize
  \bibliography{../Bibliography_base}
\end{frame}

%% ============================================================
\end{document}
